% =========================================================
% Compléments M1 — Calibration (C. Dérivations)
% =========================================================

\subsection*{(1) Moindres carrés pondérés : gradient et Hessien (lecture Gauss--Newton)}
Considérons une perte de type moindres carrés pondérés :
\[
\mathcal{L}(\theta)=\frac12\sum_{i=1}^N w_i\,r_i(\theta)^2
=\frac12\,r(\theta)^\top W\,r(\theta),
\]
où $r_i(\theta)$ est un résidu (prix ou IV) et $W=\mathrm{diag}(w_i)$.
Notons $J(\theta)$ le jacobien $J_{ij}=\partial r_i/\partial \theta_j$.

\paragraph{Gradient.}
En dérivant :
\[
\nabla \mathcal{L}(\theta)=J(\theta)^\top W\,r(\theta).
\]
\paragraph{Hessien exact.}
Il s'écrit
\[
\nabla^2\mathcal{L}(\theta)=J^\top WJ + \sum_{i=1}^N w_i r_i(\theta)\,\nabla^2 r_i(\theta).
\]
\paragraph{Approximation Gauss--Newton.}
Près d'un bon fit, les résidus $r_i$ sont petits, donc on néglige le second terme :
\[
\nabla^2\mathcal{L}(\theta)\approx J^\top WJ.
\]
Cela justifie beaucoup d'algorithmes pratiques : on approxime la courbure par $J^\top WJ$.

\subsection*{(2) Levenberg--Marquardt : compromis entre Newton et descente de gradient}
Gauss--Newton propose de résoudre
\[
(J^\top WJ)\Delta\theta = -J^\top Wr.
\]
Si $J^\top WJ$ est mal conditionné, on ajoute un terme de régularisation (damping) :
\[
(J^\top WJ+\lambda I)\Delta\theta = -J^\top Wr,\qquad \lambda>0.
\]
\paragraph{Interprétation.}
\begin{itemize}[leftmargin=*]
\item $\lambda$ petit : comportement proche de Newton/Gauss--Newton (rapide si bien conditionné).
\item $\lambda$ grand : comportement proche d'une descente de gradient (robuste mais lent).
\end{itemize}
En calibration, ce mécanisme est crucial : les surfaces sont bruitées et les pricers peuvent être irréguliers, donc on a besoin de robustesse.

\subsection*{(3) Log-paramétrisation : rendre les contraintes «naturelles»}
Beaucoup de paramètres doivent être positifs : $\kappa,\theta,\sigma,v_0,\alpha,\nu,\dots$.
Plutôt que d'imposer des bornes dures, on peut optimiser sur des variables libres :
\[
\kappa=\exp(\widetilde\kappa),\quad \theta=\exp(\widetilde\theta),\quad \sigma=\exp(\widetilde\sigma),\ \dots
\]
Cette reparamétrisation :
\begin{itemize}[leftmargin=*]
\item garantit la positivité ;
\item homogénéise souvent les échelles ;
\item améliore le conditionnement (car les paramètres sont naturellement en échelle log en finance).
\end{itemize}

\subsection*{(4) Mini-exercice corrigé : «fit en vol» = «fit en prix pondéré» (au 1er ordre)}
Supposons que le marché soit donné en IV et qu'on compare des IV.
Montrer qu'au premier ordre, minimiser $\sum (\sigma^{\mathrm{mod}}-\sigma^{\mathrm{mkt}})^2$ revient à minimiser une somme de carrés de prix \emph{pondérée} par $1/\mathrm{Vega}^2$.

\emph{Solution.}
Sur un point, on a la linéarisation BS :
\[
C(\sigma^{\mathrm{mod}})\approx C(\sigma^{\mathrm{mkt}})+\mathrm{Vega}(\sigma^{\mathrm{mkt}})\cdot(\sigma^{\mathrm{mod}}-\sigma^{\mathrm{mkt}}).
\]
Donc
\[
\Delta C \approx \mathrm{Vega}\cdot \Delta\sigma
\quad\Rightarrow\quad
\Delta\sigma \approx \frac{\Delta C}{\mathrm{Vega}}.
\]
En sommant sur les points :
\[
\sum (\Delta\sigma)^2 \approx \sum \left(\frac{\Delta C}{\mathrm{Vega}}\right)^2,
\]
soit un fit en prix avec poids $w\propto 1/\mathrm{Vega}^2$.
Cela explique pourquoi «fit en vol non pondéré» sur-pondère mécaniquement les zones où Vega est faible.
